\section{Introduction}
\label{sec:introduction}

Power-flow feasibility asks whether the electrical laws and operating limits
of a network can hold simultaneously. Even without switching decisions,
power is nonlinear in voltage. This leaves two separate questions: how hard
is it to decide exact feasibility, and how much does a small numerical
residual establish? We answer both for a resistive model with rational
conductances, positive voltage intervals, and signed injection intervals,
and transfer the exact classification to several precisely specified AC
models.

Our main result is that exact resistive feasibility is complete for
$\ER$, the class associated with the existential theory of the reals.
For each fixed integer $r$, hardness holds simultaneously on connected
simple planar bipartite graphs of maximum degree three and girth at least
$r$, with unit conductance and interval endpoints from fixed finite rational
sets depending only on $r$ (Theorem~\ref{thm:structural}). Thus neither
large numerical coefficients nor dense graphs are required. The class
satisfies $\NP\subseteq\ER\subseteq\PSPACE$; its relation to exact
real-valued certificates and its range of applications are surveyed by
\citet{SchaeferEtAl2024}. In particular, a polynomial-size certificate scheme
for every feasible instance, with deterministic polynomial-time verification,
would imply $\NP=\ER$.

\subsection{Results and proof strategy}

The construction starts from bounded addition and inversion. A voltage-pinned
bus imposes an affine relation on neighboring voltages. Alternating complement
paths supply distinct copies of a source variable without increasing the
degree beyond three. One small nonlinear gadget enforces inversion. Every
source solution has exactly one network extension, and the original variables
are read directly from designated bus voltages
(Section~\ref{sec:resistive}). A bounded arithmetic crossover gives planarity;
free-injection connectors give connectedness; and harmonic subdivisions give
unit conductance, bipartiteness, and prescribed girth
(Section~\ref{sec:structural}). These transformations preserve the rational
solution-set correspondence. The hard graphs may have many cycles: the girth
statement does not assert hardness on trees.

For AC feasibility, a local cosine inequality constrains the principal angle
of each line but need not provide consistent real bus angles. A cycle can
wind once around the origin while every local inequality holds. We give a
polynomial-size existential-real encoding of the missing cycle condition,
including unequal voltage magnitudes and every rational cosine in $(-1,1]$.
For purely resistive lines, a real-angle energy identity then forces equal
angles when all reactive injections are zero. This transfers the resistive
classification to AC feasibility. The same transfer allows all reactive
intervals to have positive width on a common side of zero. We also treat a
rectangular bus-angle box and principal-angle windows that are strictly
positive and shrink with the number of buses
(Section~\ref{sec:ac}). The latter windows have polynomial-bit rational
cosines; unlike the fixed-data resistive construction, those cosines need
not belong to a fixed finite set.

The solution-set correspondence yields a sharp algebraic statement
(Section~\ref{sec:algebraic}). A compact semialgebraic set defined over
$\Q$ is rationally equivalent to a feasible voltage set if and only if
it is basic closed over $\Q$, meaning that a conjunction of polynomial
equalities and weak inequalities suffices to describe it. Arbitrary compact
semialgebraic sets still occur up to semialgebraic homeomorphism. Distinguishing
these two forms of universality is necessary: rational equivalence preserves
more than topology. In particular, a rational network may have a unique
feasible voltage vector with a designated voltage generating any prescribed
real algebraic number field. Appendix~\ref{app:arithmetic} supplies the
conjunction-only arithmetic construction used for these conclusions.

Finally, Section~\ref{sec:numerical} quantifies the difference between exact
and approximate feasibility. The ordinary reduction has a two-sided residual
bound. An explicit infeasible family of $O(k)$ buses admits states with
residual less than $2^{-2^k}$. A general polynomial-minimum bound gives the
matching doubly exponential separation scale for fixed-data bounded-degree
networks. A small residual therefore cannot serve as a universal exact
feasibility test at polynomial precision. At a specified gap promise,
however, short rational certificates exist. A quantitative energy estimate
also controls angle deviation and active-power error when reactive injections
are small. Appendix~\ref{app:verification} records exact-arithmetic checks
and analytically resolves the examples used in earlier numerical experiments.

\subsection{Relation to previous work and model scope}

The resistive equations are the physical direct-current power-flow model
studied by \citet{GanLow2014}. They are distinct from the linear ``DC
approximation'' of AC power flow. Gan and Low give sufficient conditions for
exactness of a second-order cone relaxation of resistive optimal power flow,
including conditions involving voltage upper bounds and injection lower
bounds. Our classification concerns arbitrary allowed bus intervals and does
not assume those exactness conditions. Both voltage and injection bounds may
be singletons, independently; in particular, a voltage-pinned bus may also
have a prescribed injection. Section~\ref{sec:foundations} makes these input
choices explicit because the reduction uses them.

A different resistive model fixes source voltages and asks which
constant-power demands can be supplied at load buses.
\citet[Theorem 3.22]{JeeningaEtAl2023} characterize both exact feasibility and
its interior by matrix-inequality alternatives, and establish convexity of
the feasible demand set. This is a different feasible set from the one
obtained by placing voltage and injection intervals at every bus type.
Nor does an exact matrix-inequality characterization alone establish a
polynomial-time Turing algorithm for exact feasibility.

For AC networks, \citet{BienstockVerma2019} prove strong NP-hardness in a
lossless model with fixed voltage magnitudes and unconstrained reactive
power. \citet{LehmannEtAl2016} prove NP-hardness on trees, already using star
networks, in a fixed-magnitude model with active and reactive constraints.
These results concern different physical restrictions from our resistive-line
AC subclass, where variable magnitudes implement arithmetic and reactive
balance forces equal angles. Our rational-cosine encoding also states
explicitly how angle limits are represented in a binary input.

The connection between resistive AC networks and zero reactive injections
appears in \citet[Appendix B, Case 2]{LavaeiLow2012}. The reactive equations
are also weighted sine-coupling equilibrium equations, as in the oscillator
framework of \citet{DorflerEtAl2013}. We use a self-contained energy proof
with a consistent real-angle lift. The distinction is essential on meshed
networks: the winding examples in Section~\ref{subsec:principal} satisfy
zero reactive injection with nonconstant phasors.

Approximation results address a further distinction.
\citet[Theorem 7 and Corollary 8 of the arXiv version]{BienstockMunoz2018}
give LP approximation
schemes for polynomial network optimization and bounded-treewidth AC optimal
power flow, with scaled feasibility and optimality tolerances. Their
guarantees do not decide exact feasibility. Our gap-promise certificate
result likewise states its tolerance and promise explicitly; it does not
resolve the unpromised approximate-membership question posed for the
lossless fixed-magnitude model in
\citet[Section 1.3]{BienstockVermaPreprint2019}.

The arithmetic starting problem comes from
\citet{AbrahamsenEtAl2018,AbrahamsenEtAl2022}; its planar crossover comes from
\citet{DobbinsEtAl2023}. For solution-set universality, we develop the
conjunction-only part of the arithmetic method of
\citet{AbrahamsenMiltzow2019}, with explicit gate ranges and a proof of the
basic-closed restriction. Semialgebraic triangulation
\citep{OhmotoShiota2017} supplies the separate topological statement, and
the polynomial-minimum theorem of \citet{JeronimoEtAl2013} supplies the
general residual lower bound. The electrical constructions and their
quantitative properties are proved directly below.
